\chapter{Safety and Assurance}

\section*{학습 목표}

이 장은 VLA safety를 prompt filter나 모델의 선의로 다루지 않는다. VLA가 물리 로봇을 움직이는 순간 safety는 perception, planning, control, runtime monitoring, data curation, evaluation이 함께 만드는 stack property가 된다. 안전한 VLA는 위험한 명령을 거절할 뿐 아니라, 애매하면 묻고, 불확실하면 멈추고, 실패하면 보수적으로 복구해야 한다.

\begin{enumerate}
  \item VLA safety를 semantic, geometric, dynamic, system, data safety로 나눈다.
  \item unsafe success와 safe failure를 구분한다.
  \item safety shield, runtime monitor, fallback policy를 수식으로 정의한다.
  \item assurance case와 safety reporting standard를 설계한다.
\end{enumerate}

\section{Safety is a stack property}

VLA model이 위험한 문장을 거절한다고 해서 robot이 안전해지는 것은 아니다. 반대로 model이 안전한 instruction을 받았더라도 perception error, wrong frame, stale action, controller saturation 때문에 unsafe event가 발생할 수 있다. 따라서 safety는 다음 stack 전체의 성질이다 \cite{ames2019,shielding2018,predictivesafetyfilter2021,iso10218,isots15066}.
\begin{equation}
  \mathcal{S}_{VLA}
  =
  (
  \mathcal{S}_{sem},
  \mathcal{S}_{geo},
  \mathcal{S}_{dyn},
  \mathcal{S}_{sys},
  \mathcal{S}_{data}
  ).
  \label{eq:ch19-safety-stack}
\end{equation}
각 항은 semantic safety, geometric safety, dynamic safety, system safety, data safety를 뜻한다.

\begin{table}[h]
  \centering
  \caption{VLA safety taxonomy}
  \begin{tabularx}{\linewidth}{p{0.22\linewidth}X X}
    \toprule
    Safety layer & 보는 것 & 대표 failure \\
    \midrule
    Semantic & instruction, intent, policy rule & 위험한 명령 수행, ambiguity 무시 \\
    Geometric & workspace, collision, reachability & 충돌 경로, unreachable target \\
    Dynamic & speed, force, torque, contact stability & 과도한 힘, slip, fall, unstable contact \\
    System & monitor, fallback, override, logging & watchdog 부재, emergency stop 실패 \\
    Data & demonstration, label, distribution, bias & unsafe behavior imitation, missing negative case \\
    \bottomrule
  \end{tabularx}
  \label{tab:ch19-safety-taxonomy}
\end{table}

이 taxonomy의 목적은 책임을 분리하는 것이다. semantic filter로 collision을 막을 수 없고, collision checker로 위험한 human instruction을 이해할 수 없다. 안전은 여러 gate가 서로 다른 failure mode를 막을 때 생긴다.

\section{Unsafe success and safe failure}

robot benchmark에서 success는 task outcome만 보는 경우가 많다. 그러나 safety 관점에서는 success와 safety를 분리해야 한다.
\begin{equation}
  y_i
  =
  (
  y_i^{task},
  y_i^{safe}
  ),
  \qquad
  y_i^{task}, y_i^{safe} \in \{0,1\}.
  \label{eq:ch19-task-safe-label}
\end{equation}
네 가지 경우가 있다.

\begin{table}[h]
  \centering
  \caption{Task outcome과 safety outcome의 분리}
  \begin{tabularx}{\linewidth}{p{0.18\linewidth}p{0.18\linewidth}X}
    \toprule
    Task & Safety & 해석 \\
    \midrule
    success & safe & 원하는 deployment outcome \\
    success & unsafe & unsafe success; task는 되었지만 배치하면 안 됨 \\
    failure & safe & safe failure; 멈춤, 질문, 보수적 abort \\
    failure & unsafe & 가장 나쁜 경우; task도 실패하고 위험도 발생 \\
    \bottomrule
  \end{tabularx}
  \label{tab:ch19-safe-failure}
\end{table}

Real VLA에서는 unsafe success가 특히 중요하다. robot이 컵을 빠르게 건네는 데 성공했지만 사람 손에 과도한 힘으로 밀었다면, success rate는 높아도 system은 unsafe이다.

\section{Risk model}

hazard event를 $H_j$라고 하자. 하나의 episode에서 risk는 hazard probability와 severity의 조합으로 볼 수 있다.
\begin{equation}
  R
  =
  \sum_{j=1}^{J}
  P(H_j \mid \ell,\obs_{0:T},\pi_{\theta})
  \cdot
  C(H_j).
  \label{eq:ch19-risk}
\end{equation}
$C(H_j)$는 hazard severity이다. 같은 probability라도 사람과의 접촉, sharp object, hot object, high-speed motion은 cost가 다르다.

deployment 조건은 다음과 같이 쓸 수 있다.
\begin{equation}
  \pi_{\theta}
  \in
  \Pi_{deploy}
  \quad
  \text{only if}
  \quad
  R \leq R_{max}
  \wedge
  P(H_{critical}) \leq \epsilon.
  \label{eq:ch19-deploy-risk}
\end{equation}
이 식은 risk를 완전히 정확히 계산할 수 있다는 뜻이 아니다. 최소한 어떤 hazard를 줄이려는지, 어떤 threshold를 기준으로 배치하는지 명시해야 한다는 뜻이다.

\section{Semantic safety}

semantic safety는 language instruction과 context가 허용 가능한지 판단한다. instruction $\ell$에 대해 semantic safety classifier를 다음과 같이 정의할 수 있다.
\begin{equation}
  s_{sem}
  =
  \sigma_{\phi}
  (\ell,\obs_t,m_t)
  \in
  \{
  \mathrm{allow},
  \mathrm{ask},
  \mathrm{refuse},
  \mathrm{escalate}
  \}.
  \label{eq:ch19-semantic-safety}
\end{equation}
중요한 것은 context dependence이다. ``칼을 가져와''는 주방 정리 task에서는 안전할 수 있지만, 사람을 향해 빠르게 전달하는 상황에서는 위험할 수 있다. semantic safety는 단어 금지 목록이 아니라 intent, object, user, scene, robot capability를 함께 보는 판단이다.

\begin{warningbox}{Ambiguity is a safety signal}
instruction이 애매할 때 임의로 해석해서 실행하는 것은 unsafe behavior이다. VLA가 모를 때 물어보는 능력은 대화 기능이 아니라 safety mechanism이다.
\end{warningbox}

\section{Geometric and dynamic safety}

geometric safety는 target과 path가 workspace constraint를 만족하는지 본다.
\begin{equation}
  \conf_t \in \mathcal{Q}_{free},
  \qquad
  x_{ee,t} \in \mathcal{X}_{task},
  \qquad
  d(\mathcal{B}_{robot},\mathcal{B}_{human}) \geq d_{min}.
  \label{eq:ch19-geometric-safety}
\end{equation}
dynamic safety는 속도, 힘, 토크, contact constraint를 본다.
\begin{equation}
  \|\dot{\conf}_t\| \leq v_{max},
  \qquad
  \|\tau_t\| \leq \tau_{max},
  \qquad
  \|F_{contact,t}\| \leq F_{max}.
  \label{eq:ch19-dynamic-safety}
\end{equation}
VLA action은 이런 constraint를 직접 만족하지 않을 수 있다. 따라서 raw action은 실행 전에 projection 또는 rejection을 거쳐야 한다.

\section{Safety shield}

safety shield는 learned policy가 낸 raw action을 safety constraint 안으로 넣는 layer이다.
\begin{equation}
  \act_t^{exec}
  =
  \arg\min_{\act \in \mathcal{A}_{safe}(\bel_t)}
  \|\act - \act_t^{raw}\|_W^2.
  \label{eq:ch19-safety-shield}
\end{equation}
projection이 가능하면 가장 가까운 safe action을 실행한다. projection이 불가능하면 halt, ask, fallback 중 하나를 선택한다.

control barrier function 형태로 safety set을 표현할 수도 있다. safety condition을 $h(\state)\geq0$로 두면 continuous control에서 다음 constraint를 둘 수 있다.
\begin{equation}
  \nabla h(\state) f(\state)
  +
  \nabla h(\state) g(\state)\ctrlcmd
  +
  \alpha h(\state)
  \geq 0.
  \label{eq:ch19-cbf}
\end{equation}
VLA가 직접 torque를 내지 않더라도, downstream controller는 이 constraint를 만족하는 command만 실행해야 한다.

\section{Runtime monitor}

runtime monitor는 policy 밖에서 system state를 감시한다.
\begin{equation}
  s_t^{mon}
  =
  \mathcal{M}
  (
  \bel_t,
  \act_t^{raw},
  \act_t^{exec},
  h_t,
  p_t
  ),
  \label{eq:ch19-monitor}
\end{equation}
여기서 $h_t$는 hardware state, $p_t$는 progress state이다. monitor output은 다음 mode 중 하나가 될 수 있다.
\begin{equation}
  s_t^{mon}
  \in
  \{
  \mathrm{continue},
  \mathrm{slow},
  \mathrm{replan},
  \mathrm{ask},
  \mathrm{fallback},
  \mathrm{estop}
  \}.
  \label{eq:ch19-monitor-state}
\end{equation}

\begin{table}[h]
  \centering
  \caption{Runtime monitor가 감시해야 할 신호}
  \begin{tabularx}{\linewidth}{p{0.24\linewidth}X X}
    \toprule
    Signal & 위험 의미 & 대응 \\
    \midrule
    Pose uncertainty 증가 & object 또는 robot state 불확실 & slow, re-detect, ask \\
    Action saturation & policy-controller mismatch & scale down, fallback \\
    Contact force spike & unexpected collision 또는 grasp failure & stop, retract \\
    Human proximity 감소 & 사람과 너무 가까움 & speed limit, halt \\
    Deadline miss & stale action 가능성 & reject action, hold, fallback \\
    Repeated failure & recovery loop에 빠짐 & replan, human handoff \\
    \bottomrule
  \end{tabularx}
  \label{tab:ch19-monitor-signals}
\end{table}

\section{Human supervision}

human-in-the-loop은 단순히 사람이 옆에 있다는 뜻이 아니다. human authority boundary가 정의되어야 한다.
\begin{equation}
  a_t^{final}
  =
  \mathcal{H}
  (
  \act_t^{exec},
  s_t^{mon},
  u_t^{human}
  ).
  \label{eq:ch19-human-supervision}
\end{equation}
$u_t^{human}$은 approval, correction, emergency stop, teleoperation takeover일 수 있다.

supervision design에서 중요한 질문은 세 가지다.
\begin{enumerate}
  \item 어떤 event에서 robot이 사람에게 물어보는가?
  \item 어떤 event에서 사람이 robot을 즉시 override할 수 있는가?
  \item human intervention이 data engine에 어떻게 기록되는가?
\end{enumerate}
human override가 log에 남지 않으면 safety가 실제 model capability처럼 보이는 reporting bias가 생긴다.

\section{Data safety}

VLA는 demonstration에서 behavior를 배운다. demonstration이 항상 안전하다는 보장은 없다. operator가 shortcut을 쓰거나, near-miss를 무시하거나, 안전하지 않은 speed로 task를 수행할 수 있다.
\begin{equation}
  \mathcal{D}_{safe}
  =
  \{
  \tau_i
  \in
  \mathcal{D}
  \mid
  r_{safe}(\tau_i) \geq \rho
  \}.
  \label{eq:ch19-safe-data}
\end{equation}
unsafe demonstration을 단순히 제거하는 것도 항상 답은 아니다. robot은 무엇을 하지 말아야 하는지도 배워야 한다. 따라서 data에는 positive success, safe failure, unsafe near-miss, human correction이 구분되어야 한다.

\begin{table}[h]
  \centering
  \caption{Safety-aware data label}
  \begin{tabularx}{\linewidth}{p{0.22\linewidth}X X}
    \toprule
    Label & 의미 & 학습에서의 사용 \\
    \midrule
    Safe success & 성공하고 안전함 & behavior cloning, positive example \\
    Safe failure & 실패했지만 안전하게 중단 & abstention, recovery, ask 학습 \\
    Unsafe near-miss & 사고 직전 또는 constraint violation 근접 & shield training, negative mining \\
    Human correction & 사람이 개입해 수정 & preference, recovery data \\
    Unsafe success & task는 되었지만 위험함 & success metric에서 제외, penalty data \\
    \bottomrule
  \end{tabularx}
  \label{tab:ch19-data-label}
\end{table}

\section{Red-teaming and safety benchmarks}

VLA red-teaming은 단순히 이상한 prompt를 넣어보는 것이 아니다. semantic ambiguity, adversarial instruction, physical uncertainty, perception blind spot, timing failure를 함께 만들어야 한다.

\begin{enumerate}
  \item unsafe instruction: 명시적으로 위험한 행동 요구
  \item ambiguous instruction: object, recipient, intent가 불분명한 명령
  \item context reversal: 같은 문장이 안전한 context와 위험한 context에서 다르게 작동하는지 확인
  \item perception stress: occlusion, reflective object, poor lighting
  \item dynamic stress: moving human, moving object, contact transition
  \item system stress: latency spike, dropped frame, network disconnect
\end{enumerate}

safety benchmark score는 다음 vector로 보고해야 한다.
\begin{equation}
  M_{safe}
  =
  (
  R_{unsafe},
  R_{near\ miss},
  R_{ask},
  R_{refuse},
  R_{fallback},
  T_{detect},
  T_{stop}
  ).
  \label{eq:ch19-safety-metrics}
\end{equation}
unsafe rate만 낮고 ask/refuse가 과도하면 robot은 쓸모가 줄어든다. 반대로 task success가 높아도 near-miss가 많으면 배치할 수 없다.

\section{Assurance case}

assurance는 ``안전할 것 같다''는 인상이 아니라 claim, evidence, argument의 구조이다.
\begin{equation}
  \mathcal{A}
  =
  (
  C,
  E,
  G
  ).
  \label{eq:ch19-assurance-case}
\end{equation}
$C$는 safety claim, $E$는 evidence, $G$는 evidence가 claim을 지지하는 argument이다.

\begin{table}[h]
  \centering
  \caption{VLA assurance case 예시}
  \begin{tabularx}{\linewidth}{p{0.22\linewidth}X}
    \toprule
    구성 요소 & 예시 \\
    \midrule
    Claim & robot은 사람 30 cm 이내에서 gripper 속도를 제한한다 \\
    Evidence & proximity sensor log, controller limit test, red-team episode \\
    Argument & monitor가 distance를 측정하고 shield가 speed command를 projection함 \\
    Assumption & sensor calibration이 유지되고 occlusion이 threshold 이하임 \\
    Residual risk & sensor blind spot에서는 emergency stop이 필요함 \\
    \bottomrule
  \end{tabularx}
  \label{tab:ch19-assurance}
\end{table}

VLA assurance의 어려움은 learned model이 open-ended instruction을 다룬다는 점이다. 따라서 assurance case는 model 자체를 완전히 증명하려 하기보다, 배치 범위, monitor, fallback, human override, logging을 함께 묶어야 한다.

\section{Incident report standard}

safety incident가 발생했을 때 다음 정보가 남아야 한다.
\begin{enumerate}
  \item instruction, scene snapshot, robot state, human proximity
  \item raw VLA action과 shield 이후 executed action
  \item monitor state, watchdog event, fallback 여부
  \item latency, dropped frame, hardware state
  \item contact force, collision distance, joint limit margin
  \item human intervention time과 intervention type
  \item root cause label: semantic, perception, planning, action, control, system, data
  \item 재현 가능 조건과 mitigation action
\end{enumerate}
incident log가 없으면 safety improvement는 anecdote가 된다. real deployment에서는 near-miss도 failure data로 취급해야 한다.

\begin{casebox}{Case study: unsafe success를 별도 label로 남기기}
로봇이 컵을 rack에 놓았지만 지나가는 동안 유리컵과 2 mm까지 접근했고 사람이 손을 뻗어 멈출 준비를 했다면, final success만으로는 좋은 결과처럼 보인다. Real VLA report에서는 이 episode를 unsafe success 또는 near-miss로 분리해야 한다. 그래야 safety shield와 data engine이 같은 실패를 다시 학습할 수 있다.
\end{casebox}

\chapterwork
{Control barrier functions, shielding, predictive safety filter, ISO safety standards를 VLA safety stack의 서로 다른 evidence tier로 읽어라 \cite{ames2019,shielding2018,predictivesafetyfilter2021,iso10218,isots15066}.}
{VLA가 위험한 instruction을 거절해도 robot safety가 보장되지 않는 이유는 무엇인가?}
{unsafe success, safe failure, human intervention, near-miss를 구분하는 incident report form을 작성하라.}

\section{요약}

VLA safety는 prompt engineering으로 끝나지 않는다. semantic safety는 위험한 명령과 ambiguity를 다루고, geometric safety는 collision과 reachability를 다루며, dynamic safety는 force와 stability를 제한한다. system safety는 monitor, watchdog, fallback, human override로 구현되고, data safety는 unsafe behavior가 학습되는 것을 막는다. 안전한 VLA는 모든 명령을 수행하는 robot이 아니라, 수행하면 안 되는 상황을 알고 멈출 수 있는 robot이다.

다음 장에서는 이런 safety와 real-time 제약이 humanoid와 whole-body control로 확장될 때 어떤 추가 문제가 생기는지 다룬다. humanoid VLA에서는 manipulation만이 아니라 balance, locomotion, gaze, bimanual coordination, human proximity가 하나의 action space 안에서 결합된다.
